\section{Categorical inputs and a common certificate framework}
\label{sec:inputs}
Let $N_i=(N_{ij}^k)_{j,k}$ and
\[
 B=\sum_iN_iN_{i^*},\qquad p_R(t)=\det(tI-B).
\]
For a commutative fusion ring the eigenvalues of $B$, counted with
multiplicity, are its formal codegrees. The same convention will be used
for a commutative signed Frobenius quotient. We write $d_i=\FPdim(X_i)$,
$D=\sum_i d_i^2$, and use $F:\Z(\C)\to\C$ and $I$ for forgetful and
induction whenever $\C$ is spherical.

\subsection{Arithmetic and sphericalization}
We use the following necessary conditions from
\cite[Section 2]{DP}. Formal codegrees of a complex fusion category are
$d$-numbers. For a monic minimal polynomial
$t^r+a_1t^{r-1}+\cdots+a_r$, this is the integer divisibility condition
$a_r^j\mid a_j^r$ for $1\le j\le r$.
Fusion eigenvalues are cyclotomic integers. For a based fusion subring
$S\subset R$, the quotient $\FPdim(R)/\FPdim(S)$ is an algebraic integer.
The normalized-norm and finite-prime field tests of
\cite[Lemmas 2.3--2.4]{DP} give exact certificates of the required failures.
These tests depend on polynomials, divisibility and factorizations over
finite fields, not numerical approximations to roots.

Weak integrality permits the pseudo-unitary reduction. More generally,
a ring with exactly two distinct simple Frobenius--Perron dimensions
has a pseudo-unitary Galois conjugate whenever it has a categorification;
we use \cite[Proposition 2.1]{DP}. In the spherical case a dimension
character in the FP Galois orbit gives the same reduction by
\cite[Lemma 2.2]{DP}. Every dimension character is nonzero on simples
and equal on dual pairs.

For an arbitrary categorification, pass to the canonical sphericalization
and quotient its ring by the relation $\varepsilon=-1$, where
$\dim\varepsilon=-1$. The resulting signed coefficients satisfy
\begin{equation}\label{eq:signed}
 |T_{ij}^k|\le N_{ij}^k,\qquad T_{ij}^k\equiv N_{ij}^k\pmod2.
\end{equation}
The selected lifts may be chosen to have second indicator $+1$ whenever
the underlying simple is self-dual. In that normalization the adjoint
parity condition becomes
\begin{equation}\label{eq:positive-lifts}
 T_{jj^*}^i\equiv N_{jj^*}^i\pmod4
 \quad(i=i^*).
\end{equation}
This is \cite[Lemma 7.1 and Corollary 7.2]{DP}. We impose no additional
sign gauge on the odd coefficients: its compatibility with this
normalization would need separate justification. Dual pairs are selected
compatibly. The rank-five commutativity argument of Section~\ref{sec:data}
applies because a spherical dimension character is nonzero on every
simple. Thus commutative signed-table enumeration is exhaustive here.

The spherical and signed character tests are exactly those of
\cite[Sections 2, 6--7]{DP}: nonvanishing, equality on dual pairs,
formal-codegree norm and field conditions, and the positive central-square
moment inequalities, including the central sign object when necessary.
A signed table that is a sign gauge of $R$ splits the sphericalization,
so the spherical exclusion of $R$ applies. Every other signed character
component must receive its own necessary-condition certificate. This
is a quantified test over all signed tables and all components, not a
conjecture that such tables suffice for associators.

\subsection{Low induction and exact moments}
In a pseudo-unitary categorification all central dimensions are positive.
Since the ring is commutative, the induction of the unit is
multiplicity-free. Its simple dimensions are $D/f$, one for each formal
codegree $f$, and all its twists equal one
\cite[Theorem 2.13]{Ostrik}. If $L$ is the matrix of their restrictions,
including the unit column, then
\begin{equation}\label{eq:low}
 L_{0a}=1,\quad \sum_iL_{ia}d_i=D/f_a,\quad
 (LL^{\mathsf T})_{0i}=B_{0i},\quad
 Q=B-LL^{\mathsf T}\ge0\ \text{entrywise},\quad Q\succeq0.
\end{equation}
Column bounds $L_{ia}\le B_{0i}$ and $L_{ia}^2\le B_{ii}$ make this a
finite enumeration. Equal formal-codegree columns may be put in
lexicographic order to remove only redundant permutations.
The low mass in $I(X_i)$ is
$\ell_i=\sum_j(LL^{\mathsf T})_{ij}d_j$.
The moment bounds of \cite[Theorem 3.1 and Proposition 3.2]{DP} include
\begin{equation}\label{eq:low-moment}
 (3d_i+1)\ell_i\le Dd_i(d_i+1)\quad(i=i^*),
 \qquad 3\ell_i\le Dd_i\quad(i\ne i^*).
\end{equation}
Their low-column implementation is \cite[Proposition 3.3]{DP}.
We retain them as cheap consequences of the common moment framework,
not as separate ring-specific methods.

For every positive integer $j$, the Ng--Schauenburg formula is
\begin{equation}\label{eq:NS}
 \sum_{Z\in\Irr\Z(\C)}[F(Z):X_i]\dim(Z)\theta_Z^j
   =D\nu_j(X_i).
\end{equation}
Here $\nu_1(X_i)=0$ for $i\ne0$; $\nu_2(X_i)=0$ for a nonself-dual
simple and is $\pm1$ otherwise. If a self-dual simple occurs oddly in
some $Y\otimes Y^*$, its second indicator is $+1$ by the adjoint-parity
result just cited. Formula~\eqref{eq:NS} is valid for spherical fusion
categories \cite[Theorem 4.1]{NS}; the nonnegative weights used below
require the stated pseudo-unitary reduction.

The first tests in this section are inexpensive filters. Their
consequences are not claimed to be logically independent or a minimal
axiomatization. In particular, we retain an arithmetic rejection when
it avoids constructing a much larger, logically stronger moment system.
